% Generated draft: requires mathematical and source review.
\documentclass[11pt]{article}
\usepackage[margin=25mm]{geometry}
\usepackage{fontspec}
\setmainfont{Latin Modern Roman}
\usepackage{amsmath,amssymb,mathtools,unicode-math}
\setmathfont{Latin Modern Math}
\usepackage{enumitem,needspace,xurl,hyperref}
\setlength{\parindent}{0pt}
\setlength{\parskip}{5pt}
\setlength{\emergencystretch}{3em}
\newcommand{\N}{\mathbb{N}}
\newcommand{\Z}{\mathbb{Z}}
\newcommand{\Q}{\mathbb{Q}}
\newcommand{\R}{\mathbb{R}}
\newcommand{\C}{\mathbb{C}}
\newcommand{\F}{\mathbb{F}}
\newcommand{\A}{\mathbb{A}}
\newcommand{\PP}{\mathbb{P}}
\newcommand{\E}{\mathbb{E}}
\newcommand{\eps}{\varepsilon}
\newcommand{\dd}{\,\mathrm{d}}
\begin{document}
\section*{UnsolvedMath ID: 30006423}
\textbf{Source problem code:} catalogue entry
\subsection{Definitions and mathematical statement}
Let \(G\) be a discrete group, with no countability assumption imposed by the supplied record. Its left regular representation on \(\ell^2(G)\) is defined by \(\lambda_g\delta_h=\delta_{gh}\), and its reduced group \(C^*\)-algebra is
\[
A_G=C_r^*(G)=\overline{\operatorname{span}\{\lambda_g:g\in G\}}^{\|\cdot\|}.
\]
The group \(G\) is \(C^*\)-simple if \(A_G\) has no nonzero proper closed two-sided ideals.

The supplied excerpt does not define either strict comparison or selflessness. In the questions below, \(\operatorname{SC}(A)\) and \(\operatorname{Selfless}(A)\) denote the properties intended in the cited report; these symbols are placeholders for definitions requiring source verification, rather than independently specified predicates.

For context concerning the missing comparison convention, if \(\mathcal T\) is a specified nonempty family of normalized tracial states on a unital \(C^*\)-algebra \(A\), one can formulate comparison using
\[
d_\tau(a)=\lim_{n\to\infty}(\tau\otimes\operatorname{Tr})(a^{1/n}),
\qquad a\in(A\otimes\mathcal K)_+,
\]
where \(\mathcal K\) denotes the compact operators on \(\ell^2(\N)\), and \(\operatorname{Tr}\) is the usual unnormalized trace. Cuntz subequivalence \(a\precsim b\) means that there are \(x_n\in A\otimes\mathcal K\) with \(x_n^*bx_n\to a\) in norm. A tracial comparison condition is
\[
\forall a,b\in(A\otimes\mathcal K)_+,\qquad
\bigl[d_\tau(a)<d_\tau(b)\text{ for every }\tau\in\mathcal T\bigr]
\Longrightarrow a\precsim b.
\]
This is explanatory notation only: the supplied record does not establish whether its strict comparison uses the canonical group trace \(\tau_G(a)=\langle a\delta_e,\delta_e\rangle\), all tracial states, quasitraces, or another convention. Consequently, the displayed comparison condition is not asserted to be the report's definition.

With the intended technical definitions still to be recovered, the extracted questions are:
\begin{enumerate}
\item Does every discrete group \(G\) satisfy
\[
\operatorname{SC}(A_G)\quad\Longleftrightarrow\quad G\text{ is }C^*\text{-simple}?
\]
\item Does every discrete group \(G\) satisfy
\[
\operatorname{Selfless}(A_G)
\quad\Longleftrightarrow\quad
\bigl(\operatorname{SC}(A_G)\ \text{and}\ G\text{ is }C^*\text{-simple}\bigr)?
\]
Here selflessness is written as a property of \(A_G\); the excerpt does not specify whether a distinguished trace or other additional data are part of that property.
\item Find a necessary and sufficient dynamical condition for selflessness in the reduced group \(C^*\)-algebra setting. The passage mentions Ozawa's result on extremely proximal actions, but does not specify the actions or spaces over which such a characterization is sought.
\item More generally, characterize selflessness of reduced crossed products by a dynamical condition or a condition on stabilizer subgroups. For an action \(G\curvearrowright X\), a point stabilizer is \(G_x=\{g\in G:g\cdot x=x\}\). The record does not specify whether the intended crossed products are \(C_0(X)\rtimes_r G\), \(C(X)\rtimes_r G\), or more general coefficient algebras, nor the hypotheses on the group, space, or action. Thus no particular class of crossed products or stabilizer criterion is imposed here.
\end{enumerate}
These are statement-only questions. Missing definitions and action hypotheses prevent a fully intrinsic formal statement from this excerpt alone.
\subsection{Short English statement}
For a discrete group, ask whether strict comparison of its reduced group \(C^*\)-algebra is equivalent to \(C^*\)-simplicity, and whether selflessness is equivalent to their conjunction. Also seek dynamical and stabilizer characterizations of selflessness, including for reduced crossed products. The intended definitions of strict comparison and selflessness, and the scope of the crossed-product question, require clarification from the cited sources.
\subsection{Sources}
\begin{itemize}
\item Local UnsolvedMath registry, permanent ID \(30006423\), source problem code \texttt{catalogue entry}, title: Selflessness and Strict Comparison for Reduced Group \(C^*\)-Algebras. This draft uses the supplied complete source record. No individual proposer is identified.
\item Supplied report citation: \emph{C*-Algebras (2026)}, Oberwolfach Reports, DOI \texttt{10.4171/owr/2025/35}, pp.\ 1887--1948. \url{https://doi.org/10.4171/owr/2025/35}. This is also the supplied literature source labeled ``Oberwolfach Report 35.'' The record separately gives proposed year 2025; the differing year indications are retained without resolution.
\item Supplied literature source labeled ``Strict comparison primary paper'': \url{https://doi.org/10.1007/s00222-025-01366-5}. The record supplies no authors or full title for this paper.
\item The original passage invokes Ozawa's result on extremely proximal actions as reference [30]. Its bibliographic entry and theorem statement are absent from the supplied record.
\item Literature assessment dated exactly \texttt{2026-08-22}: the supplied curation labels the problem open and states that recent work proves selflessness and strict comparison for broad classes, including hyperbolic groups, but does not establish the proposed general equivalences or general dynamical or stabilizer characterizations; it further states that the primary paper leaves higher-rank cases open. This is supplied dated curation, not an independently verified literature assessment. The hypotheses of the cited results and the meaning of the higher-rank cases are not supplied.
\item The registry labels the statement \texttt{corrected\_verified} and says that only formula delimiters, notation, grammar, and typography were repaired. These are supplied editorial labels, not independent mathematical verification by this draft.
\end{itemize}
\end{document}
